\documentclass[../DoAn.tex]{subfiles}
\begin{document}
\section{Đặt vấn đề}
\label{section:1.1}

Hoạt động tư vấn pháp luật đóng vai trò cung cấp hỗ trợ pháp lý, giải đáp các vướng mắc và hướng dẫn người dân tuân thủ quy định để bảo vệ quyền lợi hợp pháp. Hiện nay, nhu cầu tư vấn chuyên sâu về Luật Hôn nhân và gia đình đang ngày càng nhiều do các nguyên nhân như sự gia tăng của tỷ lệ ly hôn hay sự phổ biến của các quan hệ hôn nhân có yếu tố nước ngoài. Những yếu tố đó trực tiếp tạo ra nhu cầu cấp thiết đối với khả năng tra cứu nhanh chóng, chính xác hệ thống quy định pháp luật nói chung và các điều luật chuyên ngành về hôn nhân gia đình nói riêng.

Hệ thống pháp luật về hôn nhân và gia đình tại Việt Nam mang tính chất phức tạp, bao gồm nhiều phân lớp văn bản đan xen như Luật, Nghị định, Thông tư và Nghị quyết. Các hệ thống tài liệu này thường xuyên được cơ quan nhà nước cập nhật, sửa đổi, gây ra những khó khăn nhất định cho người dân trong quá trình tiếp cận và áp dụng vào thực tiễn. Đặc thù của ngôn ngữ văn bản quy phạm pháp luật thường chứa đựng nhiều thuật ngữ chuyên ngành phức tạp, làm hạn chế khả năng hiểu đúng quyền và nghĩa vụ hợp pháp của các cá nhân không có kiến thức chuyên môn về pháp luật. 

Việc tra cứu thông tin pháp lý trên các nền tảng trực tuyến hiện hành đòi hỏi người dùng phải xác định chính xác từ khóa hoặc nắm rõ số hiệu điều luật cụ thể. Các hệ thống tìm kiếm tài liệu truyền thống thường thiếu khả năng phân tích ngữ nghĩa, tạo ra rào cản lớn đối với người sử dụng thông thường. Đáng chú ý, các hệ thống văn bản pháp luật thường xuyên tồn tại mối quan hệ chồng chéo, sửa đổi, thay thế hoặc hướng dẫn lẫn nhau. Một số điều khoản có khả năng phát sinh mâu thuẫn trong các tình huống áp dụng cụ thể, khiến người dân gặp khó khăn trong việc xác định đâu là căn cứ pháp lý chính cần áp dụng, đâu là căn cứ bổ sung, hướng dẫn và đâu là các căn cứ không được áp dụng do đã hết hiệu lực hoặc không phù hợp với tình huống.

Hiện nay, một số hệ thống trí tuệ nhân tạo hỗ trợ tư vấn pháp lý đã được phát triển nhưng vẫn bộc lộ nhiều điểm hạn chế. Các hạn chế có thể kể đến như các điều khoản và căn cứ pháp lý chưa cập nhật theo các quy định mới nhất của hệ thống pháp luật, các căn cứ pháp lý được đưa ra chưa đủ để giải quyết các tình huống pháp luật hay quyền truy cập bị giới hạn trong các cơ quan, tổ chức và chưa được phổ biến đến người dân. Thực trạng này gây khó khăn cho người dùng trong việc tìm kiếm tư vấn pháp lý chính xác với bối cảnh thực tế của bản thân.

Việc giải quyết bài toán tra cứu và tư vấn pháp luật tự động mang lại những giá trị thực tiễn thiết thực cho cộng đồng. Xử lý những rào cản thông tin nêu trên giúp người dân nhanh chóng tiếp cận nguồn dữ liệu pháp lý minh bạch, chính thống và sát với hoàn cảnh cá nhân. Một cơ chế tư vấn luật pháp hiệu quả sẽ đóng vai trò bổ trợ đắc lực cho các hệ thống tư vấn công lập, góp phần giảm thiểu áp lực công việc cho đội ngũ cán bộ ngành luật. Khuôn mẫu giải quyết bài toán này không chỉ giới hạn trong phạm vi Luật Hôn nhân và gia đình mà còn sở hữu tiềm năng mở rộng sang nhiều phân hệ pháp luật chuyên ngành khác.
\section{Mục tiêu và phạm vi đề tài}
\label{section:1.2}

Hiện nay, các công cụ hỗ trợ tra cứu và tư vấn pháp luật tại Việt Nam đã có những bước phát triển nhất định với sự xuất hiện của các sản phẩm như AITracuuluat, AI luật và Chatbot của Bộ Tư pháp. Bên cạnh đó, các mô hình ngôn ngữ lớn phổ biến như Gemini cũng thường xuyên được người dùng sử dụng để giải đáp các vướng mắc pháp lý. Trong lĩnh vực chuyên biệt, các doanh nghiệp viễn thông lớn đã triển khai các hệ thống như Trợ lý ảo Viettel hay VNPT SmartBot phục vụ cho các tổ chức và cá nhân có thẩm quyền.

Qua quá trình khảo sát và đánh giá, các sản phẩm hiện có bộc lộ những ưu điểm và hạn chế khác nhau trong việc giải quyết các bài toán pháp lý. Công cụ AITracuuluat có khả năng nêu ra căn cứ pháp lý chính xác nhưng thường xuyên gặp tình trạng dữ liệu điều luật chưa được cập nhật mới nhất và phản hồi chưa thỏa đáng đối với các tình huống thực tế. Trong khi đó, các chatbot của Bộ Tư pháp đảm bảo tính cập nhật và chính xác cao nhưng lại bị giới hạn về số lượng lượt sử dụng miễn phí. Đối với các mô hình ngôn ngữ lớn như Gemini, câu trả lời đôi khi thiếu sự viện dẫn đầy đủ các căn cứ pháp lý cần thiết và gặp khó khăn trong việc cập nhật thời hạn hiệu lực của văn bản quy phạm pháp luật. Các hệ thống chuyên dụng của Viettel hay VNPT thì khả năng tiếp cận của người dân đại chúng còn hạn chế do quy định về quyền truy cập.

Tổng hợp từ các phân tích trên, các hệ thống tư vấn pháp luật tự động hiện nay vẫn tồn tại những hạn chế nhất định trong khả năng viện dẫn chính xác các điều luật làm căn cứ pháp luật để trả lời câu hỏi của người dùng. Điều này khiến người dùng khó có thể tin tưởng hoàn toàn vào sự tư vấn của hệ thống đối với các tình huống pháp lý thực tế cao.

Đề tài này hướng tới mục tiêu xây dựng một hệ thống chatbot tư vấn Luật Hôn nhân và gia đình có khả năng giải quyết các câu hỏi thuộc các chủ đề liên quan đến luật hôn nhân gia đình, bao gồm: Kết hôn, Quan hệ giữa vợ và chồng trong hôn nhân, Chấm dứt hôn nhân, Chế độ tài sản của vợ chồng, Quan hệ giữa cha mẹ và con, Quan hệ hôn nhân và gia đình có yếu tố nước ngoài. Những câu trả lời phải tham chiếu đúng các điều luật liên quan, giải thích ngắn gọn quan hệ giữa các văn bản viện dẫn và giải thích vì sao áp dụng các căn cứ pháp lý đó. Trọng tâm của nghiên cứu là thiết lập một cơ chế truy xuất thông tin đảm bảo mọi phản hổi đều dựa trên các các căn cứ pháp lý còn hiệu lực pháp luật tại thời điểm truy vấn, các căn cứ pháp lý được áp dụng theo đúng nguyên tắc hiệu lực các văn bản quy phạm pháp luật của hệ thống luật pháp Việt Nam và có thể xử lý các vấn đề chồng chéo, mâu thuẫn giữa các điều luật (nếu có).

Phần mềm được phát triển tập trung vào các chức năng chính bao gồm: hỏi đáp về nội dung Luật Hôn nhân và gia đình Việt Nam; phân loại dạng câu hỏi theo các chủ đề pháp lý nhờ các và truy xuất ngữ cảnh pháp lý từ đồ thị tri thức (Knowledge Graph) nhờ các agent; và sinh câu trả lời có trích dẫn điều khoản cụ thể. Hệ thống cung cấp khả năng hiển thị minh bạch luồng xử lý của các tác nhân (agent) và các công cụ được kích hoạt trong quá trình tìm kiếm câu trả lời. Công cụ cũng tích hợp cơ chế kiểm tra và cảnh báo về tình trạng hiệu lực, thứ bậc văn bản và các mối quan hệ sửa đổi, bổ sung của các căn cứ pháp lý được sử dụng. Việc thực hiện đề tài kỳ vọng đạt được độ chính xác trong viện dẫn pháp lý và khả năng xử lý xung đột văn bản ở mức trên 85\%, góp phần nâng cao chất lượng dịch vụ tư vấn pháp luật tự động cho cộng đồng.

\section{Định hướng giải pháp}
\label{section:1.3}

Từ việc xác định rõ những hạn chế của các hệ thống hiện tại và nhiệm vụ cần giải quyết ở phần \ref{section:1.2}, đồ án đề xuất định hướng tiếp cận sử dụng kiến trúc Agentic GraphRAG. Đây là kiến trúc sự kết hợp giữa mô hình Đồ thị tri thức (Knowledge Graph) để mô hình hóa mối quan hệ giữa các văn bản pháp lý và quy trình Agentic Workflow để điều phối luồng truy xuất. Việc lựa chọn định hướng này xuất phát từ khả năng của đồ thị tri thức trong việc biểu diễn trực quan các mối quan hệ phức tạp như sửa đổi, bổ sung hay thay thế giữa các văn bản luật, kết hợp cùng khả năng phân loại câu hỏi và định tuyến đến các bộ truy xuất (retriever) phù hợp của các tác nhân AI, nhằm lấy về context bao gồm những căn cứ pháp lý chính xác để trả lời cho câu hỏi người dùng.

Dựa trên định hướng công nghệ nêu trên, giải pháp tổng thể của đồ án là một ứng dụng chatbot trên nền tảng web tiếp nhận truy vấn bằng ngôn ngữ tự nhiên của người dùng. Khi tiếp nhận câu hỏi từ người dùng, các tác nhân AI sẽ phân loại câu hỏi thuộc nhóm chủ đề pháp lý nào để tự động sinh mã truy vấn Cypher tương ứng, lấy về các điều khoản pháp luật liên quan từ cơ sở dữ liệu đồ thị Neo4j. Hệ thống tự động kiểm tra tình trạng hiệu lực và đối chiếu quan hệ ưu tiên, giải quyết mâu thuẫn (nếu có) trong các căn cứ pháp lý trước khi đưa ngữ cảnh cho mô hình ngôn ngữ sinh ra câu trả lời cuối cùng. 

Đóng góp chính của đồ án là việc thiết kế và xây dựng thành công một đồ thị tri thức chuyên biệt mô hình hóa Luật Hôn nhân và gia đình Việt Nam năm 2014 cùng hệ thống các văn bản dưới luật liên quan. Bên cạnh đó, nghiên cứu cũng đóng góp một bộ dữ liệu đánh giá chuẩn (benchmark dataset) bao gồm tối thiểu 600 câu hỏi pháp lý. Kết quả đầu ra là một công cụ tư vấn pháp lý minh bạch về luồng suy luận của AI, mục tiêu đạt hiệu năng đánh giá RAGAS trên 85\% và tỷ lệ xử lý chính xác các xung đột pháp lý ở mức trên 80\%.

\section{Bố cục đồ án}
\label{section:1.4}

Phần còn lại của báo cáo đồ án tốt nghiệp này được tổ chức như sau.

Chương 2 trình bày tổng quan kết quả khảo sát các hệ thống chatbot tư vấn pháp luật hiện hành và thực trạng áp dụng trợ lý ảo để tư vấn luật pháp nói chung và luật hôn nhân gia đình nói riêng tại Việt Nam. Căn cứ vào các ưu nhược điểm của những công cụ đang có trên thị trường, các giải pháp hiện hành và lấy đó làm cơ sở để xác định rõ bài toán mà đồ án cần tập trung giải quyết.

Trong Chương 3, đồ án giới thiệu các cơ sở lý thuyết và nền tảng công nghệ lõi được ứng dụng để phát triển hệ thống phần mềm. Nội dung tập trung làm rõ các khái niệm về Đồ thị tri thức (Knowledge Graph), kỹ thuật Retrieval-Augmented Generation (RAG) và kiến trúc Agentic AI. Bên cạnh đó, chương này cũng đề cập đến các công cụ lưu trữ dữ liệu đồ thị như Neo4j và nền tảng giao diện Chainlit.

Chương 4 cung cấp nội dung chi tiết về quá trình thiết kế hệ thống và quá trình thử nghiệm hệ thống và phân tích các kết quả đánh giá hiệu năng. Dựa trên bộ dữ liệu kiểm chuẩn chuẩn hóa (benchmark dataset), chương này trình bày các số liệu thực nghiệm liên quan đến độ chính xác trong việc lựa chọn công cụ của tác nhân AI, khả năng truy xuất ngữ cảnh và chất lượng câu trả lời cuối cùng. Các kết quả này được đối chiếu trực tiếp với hệ thống chỉ tiêu đề ra thông qua bộ chỉ số RAGAS và các tiêu chí xử lý văn bản luật khác.

Chương 5 trình bày đóng góp cốt lõi của đồ án thông qua việc mô tả chi tiết giải pháp kỹ thuật và quá trình xây dựng hệ thống Agentic GraphRAG. Chương này mô tả việc thiết kế đồ thị knowledge graph để mô hình hóa mối quan hệ giữa Luật hôn nhân và gia đình năm 2014 và các văn bản dưới luật liên qua và cơ chế tự động kiểm tra tình trạng hiệu lực văn bản, áp dụng các văn bản theo đúng cấp bậc các văn bản pháp luật và giải quyết mâu thuẫn giữa các điều luật (nếu có). Cuối cùng, trình bày cách xây dựng tập benchmark dataset gồm 600 câu hỏi, bao phủ các dạng câu hỏi chính liên quan đến Luật hôn nhân và gia đình.

Cuối cùng, Chương 6 tổng kết lại các kết quả đã đạt được của quá trình nghiên cứu và phát triển phần mềm. Bên cạnh việc khẳng định mức độ hoàn thành các mục tiêu ban đầu, chương này cũng đưa ra những nhận định khách quan về các mặt còn hạn chế của hệ thống hiện tại, từ đó đề xuất các hướng phát triển và mở rộng tiềm năng của công cụ tư vấn pháp lý trong tương lai.

\end{document}